\chapter{输入输出}
\section{input 输入}
\begin{lstlisting}[style=pythonstyle,numbers=none]
name = input('请输入你的名字：')
print('你好', name)
\end{lstlisting}


\section{print 输出}
\codeinline{print()} 函数的详细语法格式如下
\begin{lstlisting}[style=demostyle,numbers=none]
print (value,...,sep='',end='\n',file=sys.stdout,flush=False)
\end{lstlisting}
从上面的语法格式可以看出, value 参数可以接受任意多个变量或值, 因此 \codeinline{print()} 函数完全可以输出多个值。
例如如下代码：
\begin{lstlisting}[style=pythonstyle,numbers=none]
user_name ＝ 'Charlie'
user_age = 8
#同时输出多个变量和字符串
print("读者名：",user_name,"年龄：",user_age)
\end{lstlisting}

\subsection{分隔符}
\begin{lstlisting}[style=demostyle,numbers=none]
>>> print("读者名：" ,user_name,"年龄：",user_age,sep='|')
读者名：|Charlie|年龄：|8
\end{lstlisting}

\subsection{结尾符}
\begin{lstlisting}[style=pythonstyle,numbers=none]
print('你好', '世界')
print('你好', '世界', sep='')
print('你好', '世界', sep='', end='')
print('你好', '世界', sep='', end='\n')
print('你好', '世界', sep='', end='\t')
print('你好', '世界', sep='', end='\r')
print('你好', '世界', sep='', end='\b')
\end{lstlisting}
\begin{lstlisting}[style=demostyle,numbers=none]
你好 世界
你好世界
你好世界
你好世界
你好世界
你好世界
\end{lstlisting}

\subsection{输出位置}
\codeinline{print()} 函数默认输出到屏幕。实际上，完全可以通过改变该参数让 \codeinline{print()} 函数输出到特定文件中，
\begin{lstlisting}[style=demostyle,numbers=none]
with open('test.txt', 'w') as f:
    print('Hello, world!', file=f)
\end{lstlisting}

\subsection{格式化输出}
\begin{table}[]
\begin{tabular}{|l|l|}
\hline
\multicolumn{1}{|c|}{\textbf{转换说明符}} & \multicolumn{1}{c|}{\textbf{解释}} \\ \hline
\%d、\%i                              & 转换为带符号的十进制整数                     \\ \hline
\%o                                  & 转换为带符号的八进制整数                     \\ \hline
\%x、\%X                              & 转换为带符号的十六进制整数                    \\ \hline
\%e                                  & 转化为科学计数法表示的浮点数（e 小写）             \\ \hline
\%E                                  & 转化为科学计数法表示的浮点数（E 大写）             \\ \hline
\%f、\%F                              & 转化为十进制浮点数                        \\ \hline
\%g                                  & 智能选择使用 \%f 或 \%e 格式              \\ \hline
\%G                                  & 智能选择使用 \%F 或 \%E 格式              \\ \hline
\%c                                  & 格式化字符及其 ASCII 码                  \\ \hline
\%r                                  & 使用 repr() 函数将表达式转换为字符串           \\ \hline
\%s                                  & 使用 str() 函数将表达式转换为字符串            \\ \hline
\end{tabular}
\caption{转换说明符}
\end{table}

转换说明符（Conversion Specifier）只是一个占位符，它会被后面表达式（变量、常量、数字、字符串、加减乘除等各种形式）的值代替。

\begin{lstlisting}[style=pythonstyle,numbers=none]
# 输出整数
print('%d' % 123)
# 输出浮点数
print('%f' % 3.1415926)
# 输出字符串
print('%s' % 'hello')
# 输出字符
print('%c' % 'a')
# 输出十六进制
print('%x' % 255)
# 输出八进制
print('%o' % 255)
\end{lstlisting}

\begin{lstlisting}[style=demostyle,numbers=none]    
123
3.141593
hello
a
ff
377
\end{lstlisting}

\section{文件的内建函数和方法}
\begin{lstlisting}[style=demostyle,numbers=none]
# 文件的打开 —— open()
'''
 文件的打开模式：
 open(file, mode='r', buffering=-1, encoding=None, errors=None, newline=None, closefd=True, opener=None)

 其中 
  mode 
    r 以只读方式打开文件
    w 以写入的方式打开文件，会覆盖已存在的文件
    a 以追加的方式打开文件，如果文件存在，文件指针会放在文件的结尾
    b 以二进制模式打开文件
'''
f = open('test.txt', 'r')

# 文件的读取 —— read()
print(f.read())     # 读取全部内容
print(f.readline()) # 读取一行
print(f.readlines())# 读取所有行

# 文件的关闭 —— close()
f.close()

f = open('test.txt', 'w')
# 文件的写入 —— write()
f.write('hello world')
f.close()

# 文件指针的移动 —— seek()
f = open('test.txt', 'a')
# seek(offset, whence)  
# offset: 偏移量，
# whence: 偏移的起始位置
#   0: 文件的开头
#   1: 当前位置
#   2: 文件的结尾

print('当前文件指针位置：%s' % f.tell())
print('当前读取到一个字符，字符为：', f.read(1))
print('当前文件指针位置：', f.tell())

f.seek(0) # 移动到文件的开头
print('移动到文件开头后的文件指针位置：', f.tell())

## 显示文件指针位置—— tell()
print('当前文件指针位置：%s' % f.tell())

print('执行seek(5, 1)后：')
f.seek(5, 1) # 从当前位置向后移动5个字符
print('当前文件指针位置：%s' % f.tell())

\end{lstlisting}






